% docs/UFX_LANGUAGE_REFERENCE.tex
% UFX AUTO2 Command Language -- Formal Reference
% VSEPR-SIM v5 beta9+
%
% Purpose:
%   Authoritative specification of the UFX AUTO2 command language as it exists
%   today, written as a foundation for a future universal VSEPR-SIM scripting
%   language.  Every construct documented here is implemented and testable.
%
% Audience:
%   Scientists, engineers, and future language implementers who need to
%   understand how UFX AUTO2 commands compose, what types they operate on,
%   what guarantees they make, and how the command layer maps to the
%   underlying database and generator stack.

\documentclass[11pt,letterpaper]{article}

\usepackage{geometry}
\geometry{margin=1.1in}
\usepackage{amsmath}
\usepackage{amssymb}
\usepackage{listings}
\usepackage{xcolor}
\usepackage{hyperref}
\usepackage{booktabs}
\usepackage{longtable}
\usepackage{parskip}
\usepackage{microtype}
\usepackage{titlesec}
\usepackage{enumitem}
\usepackage{fancyhdr}

\hypersetup{
  colorlinks=true,
  linkcolor=blue!70!black,
  urlcolor=blue!70!black,
  citecolor=blue!70!black
}

\definecolor{codeblue}{rgb}{0.08,0.25,0.52}
\definecolor{codegray}{rgb}{0.45,0.45,0.45}
\definecolor{codegreen}{rgb}{0.0,0.45,0.08}
\definecolor{codeback}{rgb}{0.97,0.97,0.97}
\definecolor{keywordred}{rgb}{0.65,0.08,0.08}

\lstdefinelanguage{ufxlang}{
  morekeywords={ufx,auto2,init,audit,randomfill,validate,promote,
				validate-web,fill-molecular,fill-thermo,fill-crystal,
				fill-transport,fill-mechanical,fill-macro,score,
				pipeline,block,record,axis,region,tier},
  morekeywords=[2]{--db,--batch,--count,--seed,--min-score,--warn-score,
				   --verbose,--quiet,--recompute,--web-validate,
				   --cache,--run-dir,--validate-batch,--web-batch},
  sensitive=true,
  morecomment=[l]{\#},
  morestring=[b]",
  morestring=[b]',
}

\lstset{
  language=ufxlang,
  basicstyle=\ttfamily\small,
  keywordstyle=\color{codeblue}\bfseries,
  keywordstyle=[2]\color{keywordred},
  commentstyle=\color{codegray}\itshape,
  stringstyle=\color{codegreen},
  backgroundcolor=\color{codeback},
  frame=single,
  rulecolor=\color{black!20},
  framesep=4pt,
  xleftmargin=6pt,
  xrightmargin=4pt,
  breaklines=true,
  showstringspaces=false,
  tabsize=4,
  captionpos=b,
}

\pagestyle{fancy}
\fancyhf{}
\rhead{\texttt{UFX AUTO2 Language Reference}}
\lhead{VSEPR-SIM v5 beta9+}
\cfoot{\thepage}

% ============================================================================
\begin{document}
% ============================================================================

\begin{titlepage}
\vspace*{2cm}
{\huge\bfseries UFX AUTO2 Command Language\\[0.4em]
\Large Formal Reference and Scripting Foundation\par}
\vspace{1.5cm}
{\large VSEPR-SIM Platform --- Materials State Database Generator\par}
\vspace{0.5cm}
{\normalsize Version: 5 beta9+ \quad Branch: 4.0-legacy-beta\par}
\vspace{0.3cm}
{\normalsize Document status: \textbf{Active specification}\par}
\vspace{2cm}
\hrule
\vspace{0.4cm}
{\small
This document is the authoritative reference for the \texttt{ufx auto2}
sub-language embedded inside the VSEPR-SIM executable.  It describes the
complete command grammar, data model, type system, block ontology, option
vocabulary, pipeline semantics, and correctness invariants.

It is written as a \emph{language specification} rather than a user guide,
so that future work can extend these constructs into a standalone scripting
language capable of expressing full VSEPR-SIM research workflows.
}
\vspace{0.5cm}
\hrule
\end{titlepage}

\tableofcontents
\newpage

% ============================================================================
\section{Introduction and Motivation}
% ============================================================================

\subsection{What the UFX language is}

The UFX AUTO2 command language is a deterministic, pipeline-oriented
sub-language for driving the generation, validation, enrichment, scoring,
and promotion of atomistic material state records.  It is currently embedded
in the \texttt{vsepr-ufx} executable and invoked as:

\begin{lstlisting}
vsepr ufx auto2 <command> [options]
\end{lstlisting}

The language is \emph{not} Turing-complete in its current form.  It is a
\textbf{command vocabulary} with typed options, a fixed execution model, and
a set of correctness invariants that together constitute a small but rigorous
domain-specific language (DSL) for materials database work.

\subsection{Design goals}

\begin{enumerate}[leftmargin=2em]
  \item \textbf{Deterministic.}  Given the same seed and database state,
		every command produces identical results.
  \item \textbf{Idempotent.}  Re-running any fill command against an
		already-filled database must produce zero new rows.
  \item \textbf{Composable.}  Commands chain into a well-defined pipeline
		with explicit ordering constraints.
  \item \textbf{Auditable.}  Every row written carries a \texttt{method\_tag},
		\texttt{confidence}, and provenance reference.  Nothing is hidden.
  \item \textbf{Extensible.}  The type system, block ontology, and command
		registry are designed to grow into a richer scripting language
		without breaking existing semantics.
\end{enumerate}

\subsection{Relationship to the future VSEPR-SIM scripting language}

The constructs in this document are the \emph{building blocks} of a planned
universal scripting language for VSEPR-SIM research workflows.  A future
version may expose:

\begin{itemize}
  \item First-class \texttt{pipeline} literals (ordered command sequences)
  \item Named \texttt{block} expressions with typed property access
  \item \texttt{axis} range constructors and sampler expressions
  \item \texttt{region} predicates for coverage-guided sampling
  \item Conditional execution (\texttt{if tier $\geq$ 3 then fill-thermo})
  \item Scripted promotion rules (replace the hard-coded score threshold)
\end{itemize}

Everything defined in this document has a direct, tested implementation.

% ============================================================================
\section{Command Grammar}
% ============================================================================

\subsection{Top-level invocation}

\begin{lstlisting}
vsepr ufx auto2 <command> [<option> <value>]*
\end{lstlisting}

The token sequence is:

\begin{center}
\begin{tabular}{lll}
\toprule
Position & Token & Role \\
\midrule
0 & \texttt{vsepr} & executable name \\
1 & \texttt{ufx} & subsystem qualifier \\
2 & \texttt{auto2} & version/namespace qualifier \\
3 & \texttt{<command>} & command name (see \S\ref{sec:commands}) \\
4\ldots & \texttt{<option> <value>} & key-value option pairs \\
\bottomrule
\end{tabular}
\end{center}

The qualifiers \texttt{ufx} and \texttt{auto2} are reserved for the current
namespace.  Future versions may introduce \texttt{auto3}, \texttt{viz}, or
\texttt{report} namespaces under \texttt{vsepr}.

\subsection{EBNF grammar}

\begin{lstlisting}[language={}]
invocation  ::= "vsepr" qualifier version command option*
qualifier   ::= "ufx"
version     ::= "auto2"
command     ::= "init" | "audit" | "randomfill"
			  | "validate" | "promote" | "validate-web"
			  | "fill-molecular" | "fill-thermo"  | "fill-crystal"
			  | "fill-transport" | "fill-mechanical" | "fill-macro"
			  | "score"
option      ::= flag | param
flag        ::= "--verbose" | "--quiet" | "--recompute" | "--web-validate"
param       ::= "--db"             path
			  | "--batch"          integer
			  | "--count"          integer
			  | "--seed"           uint64
			  | "--min-score"      float64
			  | "--warn-score"     float64
			  | "--validate-batch" integer
			  | "--web-batch"      integer
			  | "--cache"          path
			  | "--run-dir"        path
integer     ::= [0-9]+
uint64      ::= [0-9]+
float64     ::= [0-9]+ ("." [0-9]+)?
path        ::= <filesystem path, absolute or relative>
\end{lstlisting}

\subsection{Option parsing rules}

\begin{enumerate}[leftmargin=2em]
  \item Options are parsed left-to-right; later values override earlier ones.
  \item Unknown tokens are silently ignored at the CLI layer (forward
		compatibility).
  \item Integer overflow on \texttt{--count} or \texttt{--batch} is caught;
		the command exits with code~2.
  \item \texttt{--db} is required for all Phase~6--10 commands.  Missing
		\texttt{--db} on a fill or score command exits with code~2.
  \item Boolean flags (\texttt{--verbose}, \texttt{--recompute}) default to
		\texttt{false} and are set by presence only; there is no
		\texttt{--no-verbose} negation form in this version.
\end{enumerate}

% ============================================================================
\section{Command Reference}
\label{sec:commands}
% ============================================================================

\subsection{Overview}

\begin{longtable}{lllp{6cm}}
\toprule
Command & Phase & Tier & Description \\
\midrule
\texttt{init}            & 2   & ---  & Create or verify the database schema \\
\texttt{audit}           & 2   & ---  & Print coverage and health report \\
\texttt{randomfill}      & 2   & 0--1 & Generate candidate material records \\
\texttt{validate}        & 3--4& ---  & Run local sanity checks \\
\texttt{validate-web}    & 5   & ---  & External identity checks via PubChem/NIST \\
\texttt{promote}         & 4   & ---  & Evaluate scores and advance source\_class \\
\texttt{fill-molecular}  & 6   & 2    & Molecular descriptor block \\
\texttt{fill-thermo}     & 7   & 3    & Thermo + EOS blocks \\
\texttt{fill-crystal}    & 8   & 4    & Crystal structure block (solids only) \\
\texttt{fill-transport}  & 9   & 5    & Transport property block \\
\texttt{fill-mechanical} & 9   & 5    & Mechanical property block \\
\texttt{fill-macro}      & 9   & 5    & Alias: transport then mechanical \\
\texttt{score}           & 10  & 9    & Meta-score block and scheduler update \\
\bottomrule
\end{longtable}

\subsection{\texttt{init}}

\begin{lstlisting}
vsepr ufx auto2 init [--db <path>]
\end{lstlisting}

Creates all required tables and views in the target database using
\texttt{CREATE TABLE IF NOT EXISTS}.  Safe to re-run at any time.
Returns exit code~0 on success.

\textbf{Tables created:}
\texttt{material\_records}, \texttt{property\_values},
\texttt{validation\_records}, \texttt{generation\_axes},
\texttt{promotion\_history}, \texttt{property\_provenance}.

\textbf{Views created:} \texttt{v\_meta\_scores}.

\subsection{\texttt{audit}}

\begin{lstlisting}
vsepr ufx auto2 audit [--db <path>]
\end{lstlisting}

Queries the database and prints a structured coverage report covering:
record counts by source class, validation pass/fail rates, promotion
history summary, block fill coverage per tier, and weakest region
according to the \texttt{CoverageMap}.

\subsection{\texttt{randomfill}}

\begin{lstlisting}
vsepr ufx auto2 randomfill --count <n> [--db <path>] [--seed <uint64>]
\end{lstlisting}

Samples \texttt{count} points from the Phase~2 axis configuration
(\texttt{AxisConfig::default\_phase2()}) and inserts them as
\texttt{candidate} records into \texttt{material\_records} and
\texttt{generation\_axes}.  Deterministic when \texttt{--seed} is
provided.  Each record is assigned a unique \texttt{record\_id} derived
from its element, oxidation state, coordination, geometry, phase,
temperature, and pressure.

\subsection{\texttt{validate}}

\begin{lstlisting}
vsepr ufx auto2 validate [--batch <n>] [--db <path>] [--run-dir <path>]
\end{lstlisting}

Runs \texttt{LocalSanityChecker::check\_all()} over unvalidated records.
Writes a row to \texttt{validation\_records} for each record checked.
Operates only on records without an existing validation entry (idempotent).

\textbf{Sanity tiers checked:}
identity (Tier~0), provenance (Tier~0), state (Tier~0), force\_field
(Tier~1), molecular (Tier~2), thermo (Tier~3), crystal (Tier~4),
macro (Tier~5).

\subsection{\texttt{validate-web}}

\begin{lstlisting}
vsepr ufx auto2 validate-web [--batch <n>] [--db <path>]
							 [--cache <dir>] [--run-dir <path>]
\end{lstlisting}

Fetches external identity data from PubChem and NIST WebBook.  Results
are cached under \texttt{web\_cache/} to avoid redundant HTTP requests.
Failures are non-fatal to the pipeline (external data availability is not
a model error).

\subsection{\texttt{promote}}

\begin{lstlisting}
vsepr ufx auto2 promote [--min-score <f>] [--warn-score <f>]
						[--db <path>] [--run-dir <path>]
\end{lstlisting}

Evaluates the composite score for each validated record.  Records meeting
\texttt{min-score} transition from \texttt{candidate} to \texttt{promoted}
in \texttt{material\_records.source\_class}.  Promotion requires non-zero
provenance (hard rule: no provenance = no promotion).

\subsection{\texttt{fill-molecular} --- Phase 6}

\begin{lstlisting}
vsepr ufx auto2 fill-molecular --db <path> [--batch <n>] [--verbose]
\end{lstlisting}

Fills the \texttt{molecular} block (Tier~2) for records that do not yet
have it.  All values are computed deterministically from the periodic
table; no network calls.

\textbf{Properties written per record:}

\begin{center}
\begin{tabular}{lll}
\toprule
Property & Units & Method \\
\midrule
\texttt{molecular\_weight}      & g/mol   & atomic weight sum \\
\texttt{exact\_mass}            & Da      & monoisotopic mass \\
\texttt{heavy\_atom\_count}     & ---     & from coordination \\
\texttt{hbond\_donor\_count}    & ---     & element heuristic \\
\texttt{hbond\_acceptor\_count} & ---     & element heuristic \\
\texttt{inchikey}               & string  & stub (UFXSTUB format) \\
\texttt{canonical\_smiles}      & string  & SmilesVocab lookup \\
\bottomrule
\end{tabular}
\end{center}

\texttt{method\_tag = "atomic\_table\_phase6"}, \texttt{confidence = 0.70}.

\subsection{\texttt{fill-thermo} --- Phase 7}

\begin{lstlisting}
vsepr ufx auto2 fill-thermo --db <path> [--batch <n>] [--verbose]
\end{lstlisting}

Fills the \texttt{thermo} and \texttt{eos} blocks (Tier~3) using
Shomate polynomial heuristics and Lee-Kesler critical property
correlations.

\textbf{Properties written per record (block = \texttt{thermo}):}

\begin{center}
\begin{tabular}{lll}
\toprule
Property & Units & Method \\
\midrule
\texttt{delta\_Hf\_kJ\_mol}     & kJ/mol    & element-family seed \\
\texttt{delta\_Gf\_kJ\_mol}     & kJ/mol    & Hf + TS estimate \\
\texttt{S\_298\_J\_mol\_K}      & J/(mol·K) & Dulong-Petit proxy \\
\texttt{cp}                     & J/(mol·K) & Shomate curve, 20 pts \\
\texttt{melting\_point\_K}      & K         & family heuristic \\
\texttt{boiling\_point\_K}      & K         & family heuristic \\
\texttt{critical\_T\_K}         & K         & Lee-Kesler \\
\texttt{critical\_P\_Pa}        & Pa        & Lee-Kesler \\
\texttt{acentric\_factor}       & ---       & family estimate \\
\texttt{density\_kg\_m3}        & kg/m³     & phase estimate \\
\texttt{viscosity\_Pa\_s}       & Pa·s      & phase estimate \\
\texttt{thermal\_conductivity\_W\_mK} & W/(m·K) & Bridgman \\
\bottomrule
\end{tabular}
\end{center}

\textbf{Properties written per record (block = \texttt{eos}):}
\texttt{density\_kg\_m3} at multiple (T, P) points per the axis sample range.

\texttt{method\_tag = "shomate\_heuristic\_phase7"}, \texttt{confidence = 0.35}.

\subsection{\texttt{fill-crystal} --- Phase 8}

\begin{lstlisting}
vsepr ufx auto2 fill-crystal --db <path> [--batch <n>] [--verbose]
\end{lstlisting}

Fills the \texttt{crystal} block (Tier~4) for \emph{solid-phase records
only}.  Gas and liquid records are skipped cleanly (not an error).

\textbf{Properties written per record:}

\begin{center}
\begin{tabular}{lll}
\toprule
Property & Units & Method \\
\midrule
\texttt{space\_group}             & string  & element-family heuristic \\
\texttt{space\_group\_number}     & integer & Hermann-Mauguin map \\
\texttt{lattice\_a}               & Å       & ionic radius estimate \\
\texttt{lattice\_b}               & Å       & from symmetry \\
\texttt{lattice\_c}               & Å       & from symmetry \\
\texttt{lattice\_alpha\_deg}      & degrees & from space group \\
\texttt{unit\_cell\_volume}       & ų      & geometry product \\
\texttt{density\_g\_cm3}         & g/cm³   & mass / volume \\
\texttt{packing\_fraction}        & ---     & family seed \\
\texttt{formation\_energy\_eV\_atom} & eV/atom & element heuristic \\
\texttt{energy\_above\_hull\_eV\_atom} & eV/atom & stability estimate \\
\texttt{band\_gap\_eV}           & eV      & family heuristic \\
\texttt{debye\_temperature\_K}   & K       & Lindemann estimate \\
\bottomrule
\end{tabular}
\end{center}

\texttt{method\_tag = "crystal\_heuristic\_phase8"}, \texttt{confidence = 0.30}.

\subsection{\texttt{fill-transport} --- Phase 9}

\begin{lstlisting}
vsepr ufx auto2 fill-transport --db <path> [--batch <n>] [--verbose]
\end{lstlisting}

Fills the \texttt{transport} block (Tier~5) using Bridgman and
Kozeny-Carman correlations from prior thermo and crystal data.

\textbf{Properties written per record:}

\begin{center}
\begin{tabular}{lll}
\toprule
Property & Units & Method \\
\midrule
\texttt{k\_eff\_W\_mK}          & W/(m·K) & Maxwell-Garnett \\
\texttt{diffusivity\_m2\_s}     & m²/s    & Stokes-Einstein proxy \\
\texttt{permeability\_m2}       & m²      & Kozeny-Carman \\
\texttt{porosity}               & ---     & axis sample or default \\
\texttt{tortuosity}             & ---     & axis sample or default \\
\texttt{darcy\_coeff}           & ---     & porosity/tortuosity product \\
\texttt{mass\_transfer\_coeff}  & m/s     & film theory estimate \\
\bottomrule
\end{tabular}
\end{center}

\texttt{method\_tag = "phase9\_bridge\_heuristic"}, \texttt{confidence = 0.30}.

\subsection{\texttt{fill-mechanical} --- Phase 9}

\begin{lstlisting}
vsepr ufx auto2 fill-mechanical --db <path> [--batch <n>] [--verbose]
\end{lstlisting}

Fills the \texttt{mechanical} block (Tier~5) using UFF stiffness-to-modulus
bridge estimates and Voigt/Reuss bounds.

\textbf{Properties written per record:}

\begin{center}
\begin{tabular}{lll}
\toprule
Property & Units & Method \\
\midrule
\texttt{E\_eff\_GPa}    & GPa & Voigt bound + porosity correction \\
\texttt{G\_eff\_GPa}    & GPa & from E and Poisson ratio \\
\texttt{K\_eff\_GPa}    & GPa & bulk modulus estimate \\
\texttt{poisson\_ratio} & --- & family heuristic \\
\texttt{yield\_MPa}     & MPa & $0.6 \times E / 300$ proxy \\
\bottomrule
\end{tabular}
\end{center}

\texttt{method\_tag = "phase9\_bridge\_heuristic"}, \texttt{confidence = 0.30}.

\subsection{\texttt{fill-macro} --- Phase 9 alias}

\begin{lstlisting}
vsepr ufx auto2 fill-macro --db <path> [--batch <n>] [--verbose]
\end{lstlisting}

Executes \texttt{fill-transport} followed by \texttt{fill-mechanical} in
order.  Stops on first failure.  Transport data feeds mechanical
estimates, so order is mandatory.

\subsection{\texttt{score} --- Phase 10}

\begin{lstlisting}
vsepr ufx auto2 score --db <path> [--batch <n>] [--verbose] [--recompute]
\end{lstlisting}

Computes the \texttt{meta} block (Tier~9) for fully-filled records.
Updates the \texttt{CoverageMap} and steers the Phase~2 sampler
toward under-covered regions.

\textbf{Composite score model:}
\begin{equation}
  S_{\text{composite}} = 0.25 S_{\text{identity}} + 0.20 S_{\text{local}}
  + 0.20 S_{\text{reference}} + 0.15 S_{\text{web}}
  + 0.10 S_{\text{repeat}} + 0.10 S_{\text{macro}}
\end{equation}

\textbf{Properties written per record:}

\begin{center}
\begin{tabular}{ll}
\toprule
Property & Description \\
\midrule
\texttt{weirdness}           & Distance from nearest neighbour in axis space \\
\texttt{coverage}            & Regional record density \\
\texttt{confidence}          & Mean confidence over all filled blocks \\
\texttt{novelty}             & Inverse of hit count for this region \\
\texttt{local\_consistency}  & Variance in composite scores within region \\
\texttt{source\_conflict}    & Flag if local and web checks disagree \\
\texttt{macro\_usefulness}   & Mechanical/transport property utility score \\
\texttt{simulation\_stability} & UFF convergence quality \\
\texttt{extrapolation\_risk} & Fraction of properties below confidence threshold \\
\texttt{composite\_score}    & Weighted sum (see equation above) \\
\bottomrule
\end{tabular}
\end{center}

\textbf{Sampling priority function:}
\begin{equation}
  p(r) = (1 - c_r)^{2} \cdot (1 + \bar{u}_r) \cdot (1 - e_r) \cdot (1 + 0.2\,\bar{w}_r)
\end{equation}
where $c_r$ = coverage, $\bar{u}_r$ = mean macro usefulness, $e_r$ = extrapolation
risk, $\bar{w}_r$ = mean weirdness for region $r$.

\texttt{method\_tag = "meta\_score\_phase10"}, \texttt{confidence = 0.50}.

% ============================================================================
\section{Type System}
% ============================================================================

\subsection{Primitive types}

\begin{center}
\begin{tabular}{lll}
\toprule
UFX type & C++ type & Notes \\
\midrule
\texttt{integer}   & \texttt{int}                    & signed 32-bit \\
\texttt{uint64}    & \texttt{uint64\_t}              & seed values \\
\texttt{float64}   & \texttt{double}                 & all physical quantities \\
\texttt{string}    & \texttt{std::string}            & UTF-8 \\
\texttt{path}      & \texttt{std::filesystem::path}  & OS-native path \\
\texttt{bool}      & \texttt{bool}                   & flags only \\
\bottomrule
\end{tabular}
\end{center}

\subsection{Structured types}

\textbf{\texttt{Auto2Options}} --- the parsed option bag passed to every
Phase~6--10 handler:

\begin{lstlisting}[language=C++]
struct Auto2Options {
	std::filesystem::path db_path;
	int         batch          = 500;
	int         validate_batch = 100;
	int         web_batch      = 25;
	double      min_score      = 0.92;
	bool        verbose        = true;
	bool        recompute      = false;
	bool        web_validate   = false;
	std::string cache_dir      = "web_cache";
	std::string run_dir;
};
\end{lstlisting}

Raw subcommand strings must not pass beyond the CLI layer.  All
string-to-value conversions happen in \texttt{parse\_auto2\_options\_()}.

\textbf{\texttt{AxisSample}} --- one drawn point from the generation space:

\begin{lstlisting}[language=C++]
struct AxisSample {
	// Phase 2 (mandatory)
	std::string element_family;
	std::string element;
	int         oxidation_state;
	int         coordination;
	std::string geometry;
	std::string phase;
	double      temperature_K;
	double      pressure_atm;
	std::string target_property;
	uint64_t    seed_used;
	int         sample_index;
	// Phase 6-9 (optional, populated by generators)
	std::optional<std::string> canonical_smiles;
	std::optional<int>         heavy_atom_count;
	// ... (see axis_config.hpp for complete definition)
};
\end{lstlisting}

\textbf{Phase result types} --- each fill command returns a typed result:

\begin{lstlisting}[language=C++]
struct FillMolecularResult {
	std::string db_path;
	int processed, filled, skipped, failed;
	bool        success;
	std::string error_message;
};
// Analogous: FillThermoResult, FillCrystalResult,
//            FillMacroResult, ScoreResult
\end{lstlisting}

\textbf{Exit codes:}

\begin{center}
\begin{tabular}{ll}
\toprule
Code & Meaning \\
\midrule
0 & Success (including NoOp / already filled) \\
1 & Generator or database error \\
2 & Invalid invocation (missing required option) \\
\bottomrule
\end{tabular}
\end{center}

% ============================================================================
\section{Block Ontology}
% ============================================================================

\subsection{What a block is}

A \emph{block} is a named, versioned group of properties written to
\texttt{property\_values} for a given material record.  Every row in
\texttt{property\_values} carries:

\begin{center}
\begin{tabular}{ll}
\toprule
Column & Meaning \\
\midrule
\texttt{material\_id}   & FK to \texttt{material\_records.id} \\
\texttt{block\_tier}    & Integer tier (0--9) \\
\texttt{block\_name}    & Block identifier string \\
\texttt{property\_name} & Property identifier within the block \\
\texttt{value}          & Float64 value \\
\texttt{value\_string}  & String value (when not numeric) \\
\texttt{temperature\_K} & State temperature \\
\texttt{pressure\_atm}  & State pressure \\
\texttt{phase\_label}   & Phase label at measurement state \\
\texttt{method\_tag}    & Algorithm used to produce this value \\
\texttt{confidence}     & Estimated reliability $\in [0, 1]$ \\
\bottomrule
\end{tabular}
\end{center}

\subsection{Block name constants}

Block names are centrally defined in \texttt{ufx\_blocks} and must
not be written as raw string literals at call sites:

\begin{lstlisting}[language=C++]
namespace ufx_blocks {
	constexpr std::string_view molecular  = "molecular";
	constexpr std::string_view thermo     = "thermo";
	constexpr std::string_view crystal    = "crystal";
	constexpr std::string_view transport  = "transport";
	constexpr std::string_view mechanical = "mechanical";
	constexpr std::string_view meta       = "meta";
}
\end{lstlisting}

\textbf{Invariant:} A trailing space in a block name is a data error.
SQLite will store it silently.  Always use the constants above.

\subsection{Tier definitions}

\begin{center}
\begin{tabular}{cll}
\toprule
Tier & Block(s) & Source \\
\midrule
0 & identity, provenance, state & \texttt{randomfill} \\
1 & force\_field                 & \texttt{randomfill} (UFF) \\
2 & molecular                    & \texttt{fill-molecular} \\
3 & thermo, eos                  & \texttt{fill-thermo} \\
4 & crystal                      & \texttt{fill-crystal} \\
5 & transport, mechanical        & \texttt{fill-transport/mechanical} \\
6--8 & (reserved)               & future phases \\
9 & meta                         & \texttt{score} \\
\bottomrule
\end{tabular}
\end{center}

Each fill command targets exactly one tier.  A record missing a lower
tier does not block a higher-tier fill (fills are independent), but
lower-tier data enriches higher-tier estimates.

% ============================================================================
\section{Pipeline Semantics}
% ============================================================================

\subsection{Canonical pipeline order}

\begin{lstlisting}
randomfill
fill-molecular
fill-thermo
fill-crystal
fill-transport
fill-mechanical
validate
validate-web        # optional
score
promote
audit
\end{lstlisting}

\textbf{Ordering constraint:} \texttt{score} must run after \texttt{validate}
and \texttt{validate-web}.  Scoring before validation produces scores that
do not account for validation failures.

\textbf{Independence:} Fill steps for different tiers are independent.
Running them out of phase-order is permitted but produces lower-quality
estimates (e.g.\ crystal estimates are better when thermo density is
already filled).

\subsection{Idempotency invariant}

Every fill command implements:

\begin{lstlisting}[language=SQL]
SELECT id FROM material_records mr
WHERE NOT EXISTS (
	SELECT 1 FROM property_values pv
	WHERE pv.material_id = mr.id
	  AND pv.block_name  = ?    -- block name constant
)
LIMIT ?;
\end{lstlisting}

Re-running any fill command on an already-filled database produces exactly
zero new rows.  The result status is \texttt{NoOp}, exit code~0.

\subsection{Continual loop}

The PowerShell harness \texttt{tools/run\_ufx\_auto2\_continual.ps1}
wraps the canonical pipeline in a loop that repeats until the database
reaches a target size.  Each cycle produces:

\begin{itemize}
  \item A timestamped run directory with per-command log files
  \item A row appended to \texttt{summary.csv} with per-step exit codes
  \item An optional snapshot copy of the database
\end{itemize}

The harness checks \texttt{\$LASTEXITCODE} after each step and throws
on failure.  \texttt{validate-web} failure is non-fatal (external
network availability is not a model error).

% ============================================================================
\section{Correctness Invariants}
% ============================================================================

These invariants are enforced by the implementation and must be preserved
by any future extension.

\begin{enumerate}[leftmargin=2em]
  \item \textbf{No provenance = no promotion.}  A record without at least
		one \texttt{property\_provenance} row cannot be promoted regardless
		of its composite score.

  \item \textbf{State on every row.}  Every \texttt{property\_values} row
		must carry non-null \texttt{temperature\_K}, \texttt{pressure\_atm},
		and \texttt{phase\_label}.

  \item \textbf{Fill idempotency.}  See \S5.2.  Duplicate rows are a
		data error; the idempotency SELECT must be present in every fill loop.

  \item \textbf{Crystal only for solids.}  \texttt{fill-crystal} must
		skip gas and liquid records cleanly.  It must not emit
		\texttt{crystal\_block\_invalid} for non-solid records.

  \item \textbf{fill-macro order.}  Transport must run before mechanical.
		Transport data (density, viscosity, conductivity) feeds the
		mechanical estimate bridge.

  \item \textbf{Score after validate.}  \texttt{score} reads validation
		outcomes from \texttt{validation\_records}.  Running before
		validate produces stale scores.

  \item \textbf{Block names from constants.}  Raw block name string literals
		must not appear at call sites.  Use \texttt{ufx\_blocks::*} constants.

  \item \textbf{Phase 10 steers, not suppresses.}  The meta-score engine
		adjusts sampling priority.  It must never delete or reject records
		based on score alone.

  \item \textbf{Determinism.}  Given the same seed, \texttt{randomfill}
		produces the same records.  Given the same database state, all
		fill commands produce the same rows.  No random or time-dependent
		behavior may appear in the fill path.
\end{enumerate}

% ============================================================================
\section{Axis Space}
% ============================================================================

\subsection{Phase 2 generation axes}

The base sampling space is defined by \texttt{AxisConfig::default\_phase2()}.

\begin{center}
\begin{longtable}{llp{6.5cm}}
\toprule
Axis & Type & Values / Range \\
\midrule
\texttt{element\_family}  & discrete  & alkali, alkaline\_earth, transition\_metal,
										post\_transition, metalloid, nonmetal,
										halogen, noble\_gas, lanthanide, actinide \\
\texttt{element}          & discrete  & 56 elements from H through Pu \\
\texttt{geometry}         & discrete  & linear, trigonal\_planar, tetrahedral,
										square\_planar, trigonal\_bipyramidal,
										octahedral, pentagonal\_bipyramidal \\
\texttt{phase}            & discrete  & gas, liquid, solid, powder, molten\_salt \\
\texttt{target\_property} & discrete  & force\_field, thermo, crystal, transport, mechanical \\
\texttt{oxidation\_state} & integer   & $[-4, 8]$ \\
\texttt{coordination}     & integer   & $[1, 12]$ \\
\texttt{temperature\_K}   & float64   & $[50, 2500]$ \\
\texttt{pressure\_atm}    & float64   & $[0.001, 1000]$ \\
\bottomrule
\end{longtable}
\end{center}

\subsection{Phase 6--9 extended axes}

Additional fields on \texttt{AxisSample} are populated by the generators
and may be used to bias future sampling runs:

\begin{center}
\begin{tabular}{llll}
\toprule
Field & Phase & Type & Units \\
\midrule
\texttt{canonical\_smiles}         & 6 & string  & --- \\
\texttt{heavy\_atom\_count}        & 6 & integer & --- \\
\texttt{delta\_Hf\_kJ\_mol}        & 7 & float64 & kJ/mol \\
\texttt{cp\_298\_J\_mol\_K}        & 7 & float64 & J/(mol·K) \\
\texttt{eos\_model}                & 7 & string  & --- \\
\texttt{space\_group}              & 8 & string  & --- \\
\texttt{lattice\_a\_angstrom}      & 8 & float64 & Å \\
\texttt{packing\_fraction}         & 8 & float64 & --- \\
\texttt{energy\_above\_hull\_eV\_atom} & 8 & float64 & eV/atom \\
\texttt{porosity}                  & 9 & float64 & --- \\
\texttt{tortuosity}                & 9 & float64 & --- \\
\texttt{strain\_rate}              & 9 & float64 & s$^{-1}$ \\
\texttt{test\_environment}         & 9 & string  & --- \\
\texttt{thermal\_gradient\_K\_m}   & 9 & float64 & K/m \\
\texttt{flow\_velocity\_m\_s}      & 9 & float64 & m/s \\
\bottomrule
\end{tabular}
\end{center}

% ============================================================================
\section{Coverage and Region Model}
% ============================================================================

\subsection{Region key}

A \emph{region} is identified by a \texttt{RegionKey} string constructed
from a record's element family, phase, and oxidation bracket:

\begin{lstlisting}[language=C++]
// Example: "transition_metal|solid|ox_positive"
RegionKey make_key(const std::string& family,
				   const std::string& phase,
				   int oxidation_state);

// Derive directly from a material record
RegionKey make_key_from_identity(const UFXMaterialRecord& rec);
\end{lstlisting}

\subsection{Coverage map operations}

\begin{center}
\begin{tabular}{ll}
\toprule
Method & Semantics \\
\midrule
\texttt{load\_from\_db(db, path)} & Load hit counts from \texttt{generation\_axes} \\
\texttt{record\_hit(key)}         & Increment hit count for region \\
\texttt{count\_for(key)}          & Return hit count \\
\texttt{weakest\_region()}        & Region with lowest hit count \\
\texttt{all\_regions()}           & Set of all known region keys \\
\texttt{region\_count()}          & Number of distinct regions \\
\bottomrule
\end{tabular}
\end{center}

\subsection{Phase 10 scheduler}

The meta-score engine uses the coverage map to bias \texttt{randomfill}
sampling.  The priority weight $p(r)$ is high for regions that are
under-covered, have high macro usefulness, and carry low extrapolation risk.

The scheduler does not drop records.  It only adjusts the probability
of drawing from each region in the next cycle.

% ============================================================================
\section{Schema Reference}
% ============================================================================

\subsection{Core tables}

\begin{center}
\begin{longtable}{lp{9cm}}
\toprule
Table & Purpose \\
\midrule
\texttt{material\_records}   & One row per material state.  Carries
							   \texttt{record\_id} (derived from identity),
							   \texttt{element}, \texttt{phase},
							   \texttt{source\_class}, and timestamp. \\
\texttt{property\_values}    & All generated property data.  Foreign key
							   to \texttt{material\_records}.  One row per
							   (material, block, property, T, P). \\
\texttt{validation\_records} & Local and web validation outcomes per record. \\
\texttt{generation\_axes}    & The raw \texttt{AxisSample} that produced
							   each record, for full reproducibility. \\
\texttt{promotion\_history}  & Timestamped log of every source\_class
							   transition. \\
\texttt{property\_provenance}& Links each property group to the algorithm
							   that produced it.  Required for promotion. \\
\bottomrule
\end{longtable}
\end{center}

\subsection{Views}

\texttt{v\_meta\_scores} --- joins \texttt{material\_records} with Tier~9
rows from \texttt{property\_values} to present composite score, weirdness,
novelty, and extrapolation risk in a single queryable view.

% ============================================================================
\section{Future Language Extensions}
% ============================================================================

The following constructs are \emph{not yet implemented} but are natural
next steps given the current command vocabulary.  They are documented here
as design seeds for the universal VSEPR-SIM scripting language.

\subsection{Pipeline literal}

A named, ordered sequence of commands that can be saved and replayed:

\begin{lstlisting}
pipeline phase10_full {
	randomfill  --count 500
	fill-molecular
	fill-thermo
	fill-crystal
	fill-macro
	validate
	validate-web  --optional
	score
	promote  --min-score 0.92
	audit
}
run phase10_full --db output/ufx_auto2.sqlite
\end{lstlisting}

\subsection{Block expression}

Access a named block's properties in a condition or assignment:

\begin{lstlisting}
if block(molecular).heavy_atom_count > 6 {
	fill-thermo --extended
}
\end{lstlisting}

\subsection{Axis constructor}

Build a sampling axis inline:

\begin{lstlisting}
axis temperature_K = range(200, 1200, step=50)
axis phase = ["solid", "molten_salt"]
randomfill --count 1000 --axes temperature_K, phase
\end{lstlisting}

\subsection{Region predicate}

Filter commands to a specific region:

\begin{lstlisting}
fill-crystal --where region.family == "lanthanide"
score        --where region.coverage < 0.1
\end{lstlisting}

\subsection{Conditional promotion rule}

Replace the hard-coded score threshold with a composable expression:

\begin{lstlisting}
promote --rule {
	composite_score >= 0.90
	AND block(molecular).confidence >= 0.60
	AND NOT source_conflict
}
\end{lstlisting}

\subsection{Typed variable binding}

\begin{lstlisting}
let db     = "output/ufx_auto2.sqlite"
let batch  = 500
let target = 1GB

run phase10_full --db $db --batch $batch until db_size >= $target
\end{lstlisting}

% ============================================================================
\section{Quick Reference Card}
% ============================================================================

\begin{lstlisting}
# Full Phase 10 pipeline, batch 500
vsepr ufx auto2 init            --db ufx_auto2.sqlite
vsepr ufx auto2 randomfill      --count 500   --db ufx_auto2.sqlite
vsepr ufx auto2 fill-molecular  --batch 500   --db ufx_auto2.sqlite
vsepr ufx auto2 fill-thermo     --batch 500   --db ufx_auto2.sqlite
vsepr ufx auto2 fill-crystal    --batch 500   --db ufx_auto2.sqlite
vsepr ufx auto2 fill-transport  --batch 500   --db ufx_auto2.sqlite
vsepr ufx auto2 fill-mechanical --batch 500   --db ufx_auto2.sqlite
vsepr ufx auto2 validate        --batch 200   --db ufx_auto2.sqlite
vsepr ufx auto2 validate-web    --batch 25    --db ufx_auto2.sqlite
vsepr ufx auto2 score           --batch 500   --db ufx_auto2.sqlite
vsepr ufx auto2 promote         --min-score 0.92  --db ufx_auto2.sqlite
vsepr ufx auto2 audit           --db ufx_auto2.sqlite

# Idempotency test (second run should produce 0 new rows)
vsepr ufx auto2 fill-molecular  --batch 500   --db ufx_auto2.sqlite

# Continual loop until 1 GB
.\tools\run_ufx_auto2_continual.ps1 -TargetBytes 1GB -BatchSize 500

# Audit only
vsepr ufx auto2 audit --db ufx_auto2.sqlite
\end{lstlisting}

\vspace{1cm}
\hrule
\vspace{0.4cm}
{\small
\textbf{Document version:} 1.0 \quad
\textbf{Specification status:} Active \quad
\textbf{Platform:} VSEPR-SIM v5 beta9+, branch 4.0-legacy-beta \\
All commands described in this document are implemented in
\texttt{src/cli/cmd\_ufx.cpp} and testable via the \texttt{vsepr-ufx}
executable.
}

\end{document}
